% ============================================================================
\section{TSS 的作用与结构}
\label{sec:tss-struct}
% ============================================================================

Task State Segment（TSS）是 x86 保护模式中的一个特殊数据结构，
大小为 104 字节。它的原始设计目的是支持\term{硬件任务切换}——
CPU 可以自动保存/恢复所有寄存器到 TSS 中，实现任务切换。

但现代操作系统（Linux、Windows、macOS）都不用硬件任务切换，
而是用\term{软件任务切换}（手动保存/恢复寄存器）。原因是：

\begin{itemize}
  \item 硬件任务切换速度慢（保存/恢复全部寄存器），且不灵活
  \item 软件任务切换可以只保存必要的寄存器，性能更好
  \item 软件方式更可控，便于实现 COW fork 等高级功能
\end{itemize}

但即使不用硬件任务切换，TSS 仍然\textbf{不可或缺}——
CPU 做特权级切换（Ring 3 → Ring 0）时，\textbf{必须}从 TSS 读取
\texttt{SS0:ESP0} 来获取内核栈位置。这是硬件的硬性要求，无法绕过。

\begin{designbox}
我们的策略：整个系统只维护一个 TSS 实例，只使用其中两个字段：
\begin{itemize}
  \item \texttt{ss0} = \hex{10}（内核数据段选择子）—— 固定不变
  \item \texttt{esp0} = 内核栈顶 —— 将来做进程切换时更新为当前进程的内核栈
\end{itemize}
其余 100 字节全部保持 0。
\end{designbox}

% ----------------------------------------------------------------------------
\subsection{TSS 结构体定义}
% ----------------------------------------------------------------------------

\codefile{src/include/arch/tss.h}
\begin{lstlisting}[style=cstyle]
typedef struct __attribute__((packed)) {
    uint32_t prev_tss;  /* 0x00  链接到前一个 TSS（不用） */
    uint32_t esp0;      /* 0x04  ★ Ring 0 栈指针 */
    uint32_t ss0;       /* 0x08  ★ Ring 0 栈段选择子 */
    uint32_t esp1;      /* 0x0C  Ring 1（不用） */
    uint32_t ss1;       /* 0x10 */
    uint32_t esp2;      /* 0x14  Ring 2（不用） */
    uint32_t ss2;       /* 0x18 */
    uint32_t cr3;       /* 0x1C  页目录（不用） */
    uint32_t eip;       /* 0x20 */
    uint32_t eflags;    /* 0x24 */
    uint32_t eax, ecx, edx, ebx;      /* 0x28-0x34 */
    uint32_t esp, ebp, esi, edi;      /* 0x38-0x44 */
    uint32_t es, cs, ss, ds, fs, gs;  /* 0x48-0x5C */
    uint32_t ldt;       /* 0x60  LDT 选择子（不用） */
    uint16_t trap;      /* 0x64  调试陷阱标志（不用） */
    uint16_t iomap_base;/* 0x66  I/O 位图偏移 */
} tss_t;  /* 总计 104 字节 = 0x68 */
\end{lstlisting}

其中标注 ★ 的 \texttt{esp0} 和 \texttt{ss0} 是我们唯一需要关心的两个字段。

\begin{figure}[H]
  \centering
  \begin{tikzpicture}[
    field/.style={draw, minimum height=0.8cm, anchor=west, font=\ttfamily\small},
    label/.style={font=\small, anchor=east},
    star/.style={fill=orange!20},
    dim/.style={fill=gray!10},
  ]
    \def\w{6.5}
    % Title
    \node[font=\bfseries] at (\w/2, 0.6) {TSS 结构体（104 字节）};

    % Fields (top to bottom = offset 0x00 to 0x66)
    \node[field, dim, minimum width=\w cm] (f0) at (0, -0.3) {prev\_tss};
    \node[label] at (-0.3, -0.3) {\footnotesize 0x00};

    \node[field, star, minimum width=\w cm] (f1) at (0, -1.1) {esp0 ★};
    \node[label] at (-0.3, -1.1) {\footnotesize 0x04};

    \node[field, star, minimum width=\w cm] (f2) at (0, -1.9) {ss0 ★};
    \node[label] at (-0.3, -1.9) {\footnotesize 0x08};

    \node[field, dim, minimum width=\w cm] (f3) at (0, -2.7) {esp1, ss1, esp2, ss2};
    \node[label] at (-0.3, -2.7) {\footnotesize 0x0C};

    \node[field, dim, minimum width=\w cm] (f4) at (0, -3.5) {cr3, eip, eflags};
    \node[label] at (-0.3, -3.5) {\footnotesize 0x1C};

    \node[field, dim, minimum width=\w cm] (f5) at (0, -4.3) {eax -- edi （8 个通用寄存器）};
    \node[label] at (-0.3, -4.3) {\footnotesize 0x28};

    \node[field, dim, minimum width=\w cm] (f6) at (0, -5.1) {es, cs, ss, ds, fs, gs};
    \node[label] at (-0.3, -5.1) {\footnotesize 0x48};

    \node[field, dim, minimum width=\w cm] (f7) at (0, -5.9) {ldt, trap, iomap\_base};
    \node[label] at (-0.3, -5.9) {\footnotesize 0x60};

    % Annotations — each on its own row, no overlap
    \draw[->, thick, orange!70!black] (\w + 0.3, -1.1) -- (\w + 1.3, -1.1)
      node[right, font=\small, text width=4.5cm] {CPU 在 Ring 3$\to$0 切换时\\读取此字段作为内核栈};
    \draw[->, thick, orange!70!black] (\w + 0.3, -1.9) -- (\w + 1.3, -1.9)
      node[right, font=\small] {固定为 \hex{10}（内核数据段）};

    \node[font=\small\itshape, gray, text width=4cm, anchor=north west] at (\w + 1.3, -3.2)
      {灰色字段在软件任务切换模式\\下不使用，全部保持 0};
  \end{tikzpicture}
  \caption{TSS 结构体布局——仅 \texttt{esp0} 和 \texttt{ss0} 被实际使用}
  \label{fig:tss-layout}
\end{figure}


% ----------------------------------------------------------------------------
\subsection{TSS 描述符在 GDT 中的格式}
\label{sec:tss-gdt-descriptor}
% ----------------------------------------------------------------------------

TSS 描述符是一个\term{系统段描述符}（S 位 = 0），与代码/数据段描述符（S = 1）格式不同。
关键字段：

\begin{bitfield}
  \bit{0-3}{Type = 1001b → 32-bit TSS (Available)。被 \texttt{ltr} 加载后自动变为 1011b (Busy)}
  \bit{4}{S = 0 → 系统段（非代码/数据）}
  \bit{5-6}{DPL = 0 → 只有 Ring 0 能执行 \texttt{ltr} 指令}
  \bit{7}{P = 1 → 描述符有效}
\end{bitfield}

\noindent 合并后 access byte = \texttt{0x89} = \texttt{10001001b}。

flags 字节（gran 高 4 位）：G=0（粒度为字节，TSS 通常 < 4KiB），D/B=0（TSS 描述符中无意义），L=0，AVL=0。合并后 flags = \texttt{0x00}。
